\documentclass[11pt]{article}
\usepackage[margin=1in]{geometry}
\usepackage{amsmath,amssymb,amsthm,booktabs,microtype,hyperref}
\hypersetup{hidelinks,pdfauthor={T},pdftitle={Subgroup Safety Certification After Checkpoint Selection}}
\newtheorem{theorem}{Theorem}
\title{Subgroup Safety Certification After Checkpoint Selection}
\author{T}
\date{}
\begin{document}
\maketitle
\begin{abstract}
A release cohort loses statistical force when its results guide checkpoint
selection. We combine paired sequential tests for subgroup safety with a
false-promotion bound that accounts for reuse of the cohort. The fixed
checkpoint error depends on the subgroup claim being made. The additional
error depends on maximal leakage through the selection channel. For
adaptive binary feedback with declared fidelities, a backward recursion
bounds the leakage of the final decision, accounting for early stopping
and for different histories that select the same checkpoint. Finite
examples attain the bounds and separate the costs of the subgroup claim
and checkpoint selection. For sampling without replacement from a
frozen cohort, we derive the conditional mean bound the paired test
needs. A six-case example shows that omitting it can exceed the
nominal error allowance.
\end{abstract}

\section*{Summary and scope}
\paragraph{Question.} What is the error cost of using release-cohort
feedback to select the checkpoint that the cohort will test?
\paragraph{Finding.} Maximal leakage multiplies the fixed-checkpoint
release bound. The fixed bound depends on what the release asserts.
An all-group decision avoids the extra factor used to protect
individual group declarations.
\paragraph{Evidence and scope.} The finite-channel bound
and arithmetic checks are exact. Paired tests require a conditional
mean assumption for each unsafe group. Selection accounting requires
every cohort-derived input, including timing or stop decisions if they
reveal information. The numerical examples use synthetic data.

\section{A release test that talks to training}
Suppose a team trains an assistant to respect a person's veto and tests
each proposed release on permission-sensitive cases. Returning a
detailed failure report helps the team repair a model, but a checkpoint
chosen in response is statistically dependent on the release sample.
Teams that can afford two samples can route debugging to a development
cohort and preserve a fresh release cohort. We ask what a formal guarantee
requires when a bounded feedback channel shares the release cohort.

The analysis builds on maximal leakage \cite{issa} and on information
bounds for post-selection inference \cite{rogers}. Its operational
contribution is an explicit release calibration. It distinguishes the
effect of feedback on the final chosen checkpoint from that of
feedback visible in a longer development transcript.

\section{Risk of a false promotion}
Let $S$ denote the release sample and $C$ the final selected checkpoint,
including any test configuration chosen alongside it. Assume finite
output and sample spaces. All randomness independent of $S$ may be
conditioned on first, provided the fixed-output error guarantee then
holds uniformly over that randomness. For each possible output $c$, let
$E_c$ be the set of samples on which $c$ would be falsely promoted.
Define
\[
m_c=\max_s\Pr(C=c\mid S=s),\qquad L(C\mid S)=\sum_c m_c ,
\]
with the maximum over the support of $S$. The logarithm of $L$ is
the discrete maximal leakage of this channel \cite{issa}.

\begin{theorem}
If $\Pr(S\in E_c)\le p_c$ for each fixed output $c$, then
\[
\Pr(S\in E_C)\le\min\!\left(1,\sum_c m_c p_c\right).
\]
If the bound is $p$ uniformly, the right side is at most
$\min(1,p\,L(C\mid S))$.
\end{theorem}
\begin{proof}
Partition the event by the selected output. Each summand is
\[
\sum_{s\in E_c}\Pr(S=s)\Pr(C=c\mid S=s)
 \le m_c\Pr(S\in E_c).
\]
Summing proves both statements.
\end{proof}

If $C$ is computed from transcript $T$ and independent randomness, then
$L(C\mid S)\le L(T\mid S)$. For a fixed conditional law
$K(c\mid t)$, bound
$\max_s\sum_t K(c\mid t)\Pr(T=t\mid S=s)$ by
$\sum_t K(c\mid t)\max_s\Pr(T=t\mid S=s)$ and sum over $c$.
Thus the final-checkpoint bound can improve when different feedback
histories converge to the same decision. In contrast, if checkpoint
selection also observes an uncounted score, feedback law, or runtime,
the assumed factorization through $T$ may fail.

\subsection{When the maximum channel weight is rare}
Write $K_c(s)=\Pr(C=c\mid S=s)$ and let $\mu$ be the law of $S$.
A maximum channel weight charges the same price whether it occurs
on every sample or only on a rare one. If the sample law is known,
we can account for how much mass actually lies near that maximum.

\begin{theorem}
Given a fixed-output false-pass allowance $p_c$, the sharp
contribution bound when fractional false-pass profiles are allowed is
\[
 B_c=\min_{0\le t\le1}\left\{t p_c+
              \mathbb E_\mu[(K_c(S)-t)_+]\right\}.
\]
The total selected false-pass probability is at most $\sum_c B_c$.
The same bound applies to deterministic false-pass events. It is
sharp for those events when samples are equally likely and each
$p_c$ is an integer multiple of the sample mass.
\end{theorem}
\begin{proof}
A false-pass profile $f_c(s)$ takes values in $[0,1]$ with
$\mathbb E_\mu f_c\le p_c$. A deterministic event gives an
indicator. A fractional profile uses a fresh random coin
independent of selection conditional on $S$. Its contribution is
$\mathbb E_\mu[K_c f_c]$. For any $t\in[0,1]$,
$K_c f_c\le t f_c+(K_c-t)_+$, giving the bound. To attain it,
spend the false-pass mass on the samples with largest $K_c$,
using a partial fraction only at the boundary if necessary.
Choosing $t$ at the boundary makes the pointwise bound an
identity on that allocation. Uniform samples with a whole-number
budget need no partial atom. Each column can attain its optimum
independently, since the model imposes only the separate
false-pass budgets. The selected output events are disjoint.
\end{proof}

For an uncertain law, suppose $\mathrm{TV}(\mu,\mu_0)\le\varepsilon$
and the fixed-output error budgets hold under the actual law.
For thresholds $t_c$, let $\Phi(s)=\sum_c(K_c(s)-t_c)_+$.
A valid bound is
\[
 \sum_c t_c p_c+\mathbb E_{\mu_0}\Phi(S)
      +\varepsilon\bigl(\max_s\Phi(s)-\min_s\Phi(s)\bigr).
\]
To see this, subtract the minimum of $\Phi$ before comparing
expectations. The positive mass difference has total at most
$\varepsilon$ and can increase the result by at most the stated
range times that mass. This calculation needs the same conditional
channel on the declared support and a justified bound on law drift.

For example, let $S$ be uniform on eight samples. For one of two
checkpoints, the selection probability is $7/8$ on one sample,
$1/8$ on another, and $1/2$ on the other six. The other checkpoint
gets the complementary probability. With $p_c=1/4$ for each,
the maximum-weight bound is $7/16$. Each exact column uses one
rare maximum and one ordinary sample, giving total bound $11/32$.
If the sample law can move by total variation $1/16$, setting
both thresholds to $1/2$ yields bound $47/128$, still below $7/16$.
The verifier matched 288 column optimizations against exhaustive
subatom selection and compared the sensitivity example across
57 nominal or one-unit-shift laws.

\section{Budgeting repeated randomized replies}
Suppose response $i$ gives one bit. Conditional on any reachable
history, a deterministic query about $S$ has a binary answer, and
the mechanism returns the correct bit with probability $\lambda_i$
using a fresh coin. Take $\lambda_i\in[1/2,1]$ fixed in advance.
The per-history sum of the two maximum conditional probabilities is
at most $2\lambda_i$. For a prefix $t$ and extension $y$,
\[
\begin{aligned}
\max_s\Pr(T_i=(t,y)\mid S=s)
&\le \max_s\Pr(T_{i-1}=t\mid S=s)\\
&\quad{}\cdot
\max_{s:\,\Pr(T_{i-1}=t\mid S=s)>0}
\Pr(T_i=(t,y)\mid S=s,T_{i-1}=t).
\end{aligned}
\]
A prefix impossible for every $s$ contributes zero. Summing over
extensions and iterating yields
$L(T_d\mid S)\le\prod_{i=1}^d2\lambda_i$.
Queries may adapt to earlier responses. To support early stopping,
derive the stop decision from the released history and pad the
transcript with a fixed symbol. A stop or delay that depends directly
on the hidden cohort needs its own accounting.

\section{Bounding an adaptive selection tree}
A binary tree can give a tighter bound when the protocol stops early
or several leaves select the same result. At internal node $h$,
let $\lambda_h\in[1/2,1]$ be the declared fidelity, so that the
conditional probability of taking either branch lies between
$1-\lambda_h$ and $\lambda_h$. A leaf $\ell$ determines the final
configuration $f(\ell)$. For each possible output $c$, set
$F_c(\ell)=\mathbf 1_{\{f(\ell)=c\}}$ and work backward by
\[
 F_c(h)=\lambda_h\max\{F_c(h0),F_c(h1)\}
       +(1-\lambda_h)\min\{F_c(h0),F_c(h1)\}.
\]
\begin{theorem}
With root $r$, $L(C\mid S)\le\sum_c F_c(r)$. Equality holds
for randomized response if the sample support realizes every joint
assignment of the query truth bits at the tree's internal nodes.
\end{theorem}
\begin{proof}
Fix a sample and a target output. Conditional on reaching $h$,
the probability of the target is a convex combination of its
probabilities at the two children. A child reached with
probability zero contributes zero. Backward induction bounds
the two child probabilities by their envelopes. The largest
allowed weight on the larger envelope is $\lambda_h$, proving
the stated bound. With all truth assignments available, for
each target $c$ choose a sample with the truth bit favoring
the larger child envelope at every node. Fresh response coins
then attain $F_c(r)$. The maximizing sample can vary with $c$.
\end{proof}
For example, ask for bit $x$ at fidelity $3/4$. A zero response
selects checkpoint zero immediately. A one response leads to a
perfect report of bit $y$, which selects checkpoint $y$. The
depth-based bound is $3$. With all four truth pairs possible,
each of the three terminal transcripts has maximum probability
$3/4$, giving $L(T\mid S)=9/4$. The checkpoint maxima are
$1$ and $3/4$, giving $L(C\mid S)=7/4$. All 3,180 finite
tree configurations checked agreed with forward enumeration
under full truth support. Merging outputs strictly decreased
the multiplier in 2,924 configurations. Restricting support
to even-parity truth assignments reduced it further in 48.

\section{Charging for the claim the release makes}
Suppose a fixed candidate must pass all $G$ declared groups.
If it is unsafe for even one group, a false release implies a
false pass in that particular group. With a group error bound
$a$, the fixed-candidate false-release probability is therefore
at most $a$. This uses no independence across groups.
For a claim that at least $k$ groups are safe, choose
$G-k+1$ unsafe groups when that claim is false. The other
groups cannot supply $k$ passes, so a union bound on the
chosen subset gives fixed-candidate error at most
$(G-k+1)a$. Separate declarations that any named passing
group is safe instead require control of each false declaration.

The selection theorem multiplies the fixed-candidate bound by
$L(C\mid S)$, whether the group passes use fixed samples or
valid anytime tests. Include in $S$ every input to the test
whose dependence on selection matters. Include in $C$ all
cohort-dependent configuration choices, such as the group
definitions or rule thresholds. A conditional fixed-candidate
guarantee still has to hold for any randomness conditioned out.

For a paired test, let $D_{g,i}\in\{-1,0,1\}$ be the difference
between the candidate's and baseline's binary violation indicators
on group $g$'s case $i$. Let $\delta_g<1$ be its allowed mean
increase. Under the unsafe null
$\mathbb E[D_{g,i}\mid\mathcal F_{g,i-1}]\ge\delta_g$,
choose predictable $b_{g,i}\in[0,1/(1-\delta_g)]$ and form
\[
 M_{g,t}=\prod_{i=1}^t[1+b_{g,i}(\delta_g-D_{g,i})].
\]
Each factor is nonnegative with conditional mean at most one.
Thus $M_{g,t}$ is a nonnegative supermartingale starting at one,
and the same holds for a fixed convex mixture of such processes.
Stopping at the first crossing or a finite horizon gives probability
at most $a$ of crossing $1/a$ by that horizon. Taking increasing
horizons gives the group error bound. This applies bounded-mean
betting \cite{waudby} to a simultaneous noninferiority claim
\cite{chang}. The fixed-candidate guarantee holds for overlapping
groups, provided each conditional mean assumption is justified under
the filtration its test uses.

\subsection{Testing a fixed cohort without replacement}
Suppose the $N$ paired differences for a fixed group and checkpoint
form a finite population with mean at least $\delta$ under the unsafe
null. The test sees a uniformly random order. After $t-1$ draws let
$S_{t-1}$ be their sum and let $R_t=N-t+1$. The unobserved sum is at
least $N\delta-S_{t-1}$ under the null. Hence the next difference has
conditional mean at least
\[
 \ell_t=\max\left\{-1,\frac{N\delta-S_{t-1}}{R_t}\right\}.
\]
When $\ell_t>1$, the null is already impossible even if every unseen
difference is $+1$. The same check at the census uses the observed
total directly. At all other prefixes replace $\delta$ with $\ell_t$
in the betting factor. For $\ell_t<1$ require
$0\le b_t\le 1/(1-\ell_t)$, and for $\ell_t=1$ any finite
nonnegative $b_t$ gives a nonnegative factor. Under the null the
factor has conditional mean at most one. The stopped crossing
probability is therefore at most $a$ at threshold $1/a$, with no
additional error for the impossible-null stop. This applies
without-replacement bounded-mean betting \cite{waudby} to the paired
subgroup gate. The selection multiplier still applies to the
resulting fixed-checkpoint error bound.

For a counterexample to reusing the fixed-mean factor, take three
$+1$ and three $-1$ differences, in uniform order, with $\delta=0$.
Use bets $1/2,1,1/2,1,1,0$ at successive draws. With the uncorrected
factors, wealth crosses $3/2$ on exactly three prefix patterns.
The first draw is negative, or the first three signs are $+--$,
or the first five are $+-+--$. Their probabilities sum to $7/10$,
above the nominal bound $2/3$. The experiment compared corrected
crossing probabilities from a state recursion with direct permutation
enumeration on 411 finite cases and checked 10,536 conditional factors.

The guarantee applies to the frozen frame. A sampling rule that
chooses the next case based on its unseen difference invalidates the
conditional mean calculation. When sharing a cohort across groups,
choose a filtration for each test that preserves the uniform draw
property it uses.

In a synthetic example with four independent safe groups, rare
discordance and a margin of $2/25$, the probability that every
group passes within 72 paired cases is $0.274393$ at threshold
20 and $0.034428$ at threshold 80. The comparison isolates the
cost of demanding individual declarations when the protocol needs
only an all-group release. These thresholds illustrate the
fixed-checkpoint setting. Reusing evidence to select a checkpoint
requires the separate leakage adjustment.

In a zero-violation test, an unsafe group with violation probability
at least $\epsilon$ has chance at most $(1-\epsilon)^n$ of passing
on $n$ independent cases. For a release that requires every group,
use $a=(1-\epsilon)^n$. For the stronger claim protecting
every individual group declaration, use the family bound
$Ga$. With four feedback bits, $G=3$, $\epsilon=0.05$, and
overall error allowance $0.05$, the two calculations give
different sample requirements.
\begin{center}
\begin{tabular}{lrr}
\toprule
Per-bit correctness & All-group release & Individual group family \\
\midrule
$1/2$ & $59$ & $80$ \\
$3/4$ & $91$ & $112$ \\
$1$ & $113$ & $134$ \\
\bottomrule
\end{tabular}

\smallskip
\footnotesize Minimum cases per group.
\end{center}
For a frozen frame with at least $M$ violations among $N$ cases,
a uniformly drawn subset of size $n$ misses all violations with
probability at most $\binom{N-M}{n}/\binom{N}{n}$, taking the
numerator to be zero when $n>N-M$. With $N=20$, $M=3$, leakage
multiplier $7/4$ and overall allowance $0.05$, the fixed test must
meet $1/35$. Thirteen cases give $7/228$ and fourteen give $1/57$,
so fourteen suffice. The independent-draw calculation would require
22. The difference reflects the sampling design and the finite
population claim.

The allowance covers one final selection. Promoting multiple
checkpoints requires declaring a larger error family or budgeting
each release separately.

\section{When the leakage charge is necessary}
Take $S$ uniform on the $2^d$ binary strings and one candidate for
each string. Candidate $c$ passes only when $S=c$. Every fixed
candidate has pass probability $2^{-d}$. Reveal each bit through
independent randomized response with fidelity $\lambda$. Choose
the returned string as the candidate. The pass probability is
$\lambda^d$, and the channel has $L=(2\lambda)^d$, so the uniform
bound is exact. This is a finite adversarial fixture for a release
gate, independent of the previous independent-case model.

The checker enumerates all 3,375 three-input, three-output
channels with entries in quarters and all 27 assignments of one
pass sample per candidate. All 91,125 comparisons satisfy the
bound, and equality holds in 8,199. It independently sums the
channel for all pairs of six-bit strings, finding false-promotion
probabilities $1/64$, $729/4096$, and $1$ at fidelities
$1/2$, $3/4$, and $1$, respectively. Exact arithmetic checks the
three sample-size thresholds and the failure at one fewer case.

Both multipliers can be attained in a joint fixture. Take $S=(X,U)$,
where $X$ is a uniform three-bit string and $U$ is independently
uniform on eight atoms. Select $C$ by randomized response on $X$.
For each fixed candidate $c$, let $k-1$ safe groups always pass.
Each of the remaining $G-k+1\le8$ unsafe groups passes only when
$X=c$ and $U$ lies at its distinct assigned atom. Each unsafe
group has fixed false-pass probability $1/64$. After selection,
the false-release event is $C=X$ and one of those atoms occurs.
Its probability equals
\[
 \lambda^3\frac{G-k+1}{8}
   =(2\lambda)^3\,\frac{G-k+1}{64}.
\]
The exact checker verified 45 such witnesses across group counts,
release thresholds and response fidelities. The fixture establishes
sharpness for the abstract event and channel bounds.

\section{An auditable release protocol}
Bind the cohort, violation rubric, tested slices, response
mechanism, and a maximum of four queries before the team starts
iterating. Offer coarse diagnostic bits through the measured
channel and reserve detailed failure analysis for a separate
development set. Keep terminal release flags inside the same
accounting if any subsequent checkpoint can react to them.
A prospective comparison should measure severe violations,
useful help, false holds, and reviewer time at the same declared
false-promotion allowance. The policy is valuable when a fresh
release cohort is costly and the feedback boundary is enforceable.
Distribution drift, dependent cases within a slice, and hidden
feedback invalidate different steps and should be tested separately.

\paragraph{Review status.} Independent manuscript review is pending.

\begin{thebibliography}{9}
\bibitem{issa} Ibrahim Issa, Aaron B. Wagner, and Sudeep Kamath.
\emph{An Operational Approach to Information Leakage}.
IEEE Transactions on Information Theory 66(3), 2020.
\url{https://arxiv.org/abs/1807.07878}.
\bibitem{rogers} Ryan Rogers, Aaron Roth, Adam Smith, and Om Thakkar.
\emph{Max-Information, Differential Privacy, and Post-Selection
Hypothesis Testing}. 2016. \url{https://arxiv.org/abs/1604.03924}.
\bibitem{waudby} Ian Waudby-Smith and Aaditya Ramdas.
\emph{Estimating Means of Bounded Random Variables by Betting}.
Journal of the Royal Statistical Society Series B 86(1), 2024.
\url{https://arxiv.org/abs/2010.09686}.
\bibitem{chang} Hsin-Wen Chang and Ian W. McKeague.
\emph{Nonparametric Testing for Multiple Survival Functions with
Noninferiority Margins}. Annals of Statistics 47(1), 205--232, 2019.
\url{https://www.columbia.edu/~im2131/ps/AOS_noninf-appeared.pdf}.
\end{thebibliography}
\end{document}
